% stirling-sec4on.tex
% Definitions and (corrected) proofs of C. Stirling,
% "An introduction to decidability of higher-order matching", Sections 4--6.
% HOM project working document, 2026-11-09.
\documentclass[11pt]{article}
\usepackage[margin=2.5cm]{geometry}
\usepackage[T1]{fontenc}
\usepackage{lmodern}
\usepackage{microtype}
\usepackage{amsmath,amssymb,amsthm}
\usepackage{enumerate}
\usepackage{xcolor}
\usepackage[colorlinks=true,linkcolor=blue!50!black,citecolor=blue!50!black]{hyperref}

\newcommand{\Ground}{\mathbf{0}}
\newcommand{\ord}{\mathrm{ord}}
\newcommand{\ar}{\mathrm{ar}}
\newcommand{\Sub}{\mathrm{Sub}}
\newcommand{\depth}{\mathrm{depth}}
\newcommand{\br}{\mathrm{br}}
\newcommand{\Gm}{\mathsf{G}}
\newcommand{\bnd}{\mathrel{\downarrow}}
\newcommand{\suc}[2]{#1{\cdot}#2}
\newcommand{\iv}[2]{\pi[#1,#2]}
\newcommand{\comb}{\mathrm{comb}}
\newcommand{\pth}{\mathrm{path}}
\newcommand{\CE}{C_E}
\newcommand{\sq}{\sqsubseteq}
\newcommand{\ssq}{\sqsubset}
\newcommand{\ra}[1]{\xrightarrow{#1}}
\newcommand{\args}{\mathrm{args}}
\newcommand{\env}{\mathrm{env}}

\newcommand{\Paper}{\textsf{[paper]}}
\newcommand{\Fixed}{\textsf{[corrected]}}
\newcommand{\New}{\textsf{[new]}}
\newcommand{\Sketch}{\textsf{[sketch]}}
\newcommand{\Open}{\textsf{\color{red!60!black}[open]}}
\newcommand{\Coq}{\textsf{[Coq]}}

\theoremstyle{plain}
\newtheorem{theorem}{Theorem}[section]
\newtheorem{proposition}[theorem]{Proposition}
\newtheorem{lemma}[theorem]{Lemma}
\newtheorem{corollary}[theorem]{Corollary}
\newtheorem{fact}[theorem]{Fact}
\newtheorem{claim}[theorem]{Claim}
\theoremstyle{definition}
\newtheorem{definition}[theorem]{Definition}
\newtheorem{example}[theorem]{Example}
\theoremstyle{remark}
\newtheorem{remark}[theorem]{Remark}
\newtheorem{gap}{Gap}

\title{Stirling's decidability proof for higher-order matching,\\
Sections 4--6: definitions and corrected proofs}
\author{HOM project working document}
\date{9 November 2026}

\begin{document}
\maketitle

\begin{abstract}
This document restates Sections~4--6 of C.~Stirling, \emph{An
introduction to decidability of higher-order matching} (``the paper''),
with every definition explicit.  Every result is either proved, proved
in corrected form, or marked as open.  It collects in one place the
material of \texttt{stirling-formal.tex}, \texttt{thm37-dowod.md} and
\texttt{order5-unfolding.md} (\S\S1--10), including the analysis of
non-linear heads and the generalisation of Proposition~24.
Section~5 of the paper (tree automata) is stated without proofs, since
the decidability proof does not use it.
\end{abstract}

\tableofcontents

%======================================================================
\section*{How to read this document}
%======================================================================
\addcontentsline{toc}{section}{How to read this document}

Numbers in brackets, e.g.\ ``[S, Def.~43]'', refer to the paper.  Our
numbering is independent.  Each statement carries a status tag.
\begin{description}
\item[\Paper] Stated and proved as in the paper.  The proof is
  reproduced, possibly with more detail but without change of substance.
\item[\Fixed] The paper's statement or proof is wrong or incomplete as
  written.  We give the corrected statement or proof, and a remark says
  what was changed.
\item[\New] Not in the paper.  Added here, with a proof.
\item[\Sketch] Only a proof sketch is available (in the paper or here).
  The sketch is believed correct, but it has not been checked in full.
\item[\Open] No proof is known.  The statement is used by the paper and
  is listed among the gaps in Section~\ref{sec:summary}.
\item[\Coq] Additionally machine-checked in the Coq development
  \texttt{hom/coq} (Coq~8.18).  In Coq, the results of Sections~4
  and~6.1 are taken as axioms (G1--G10, U1--U7), and everything else is
  proved from them.
\end{description}

\paragraph{Standing restrictions.} As in [S, \S\S2--6.4]:
\begin{enumerate}[(R1)]
\item there is a single base type $\Ground$;
\item the goal $u$ contains no bound variables.
\end{enumerate}
Section~\ref{sec:restrictions} summarises how the paper removes them.

%======================================================================
\section{Background from Sections 2--3}\label{sec:background}
%======================================================================

This section only fixes notation; see the paper for details.

\begin{definition}[Types, order, arity; S, Defs.~1--2]
Types are generated by $A::=\Ground\mid A\to A$.  Every non-base type
is written $(A_1,\dots,A_n,\Ground)$ for $A_1\to\cdots\to A_n\to\Ground$.
\[
\ord(\Ground)=1,\quad \ord(A_1,\dots,A_n,\Ground)=1+\max_i\ord(A_i),
\]
\[
\ar(\Ground)=0,\quad \ar(A_1,\dots,A_n,\Ground)=\max(n,\ar(A_1),\dots,\ar(A_n)).
\]
Note that the base type has order~$1$.  For example, at order $5$ the
level-$2$ variables of a solution have order $\le2$, i.e.\ type
$\Ground\to\cdots\to\Ground\to\Ground$.
\end{definition}

\begin{definition}[Long normal form, interpolation; S, Defs.~6--9]
A term of type $(A_1,\dots,A_n,\Ground)$ is in \emph{lnf} if it is
$\lambda x_1\dots x_n.\,h\,s_1\dots s_m$ with $h$ a variable or a
constant and every $s_j$ in lnf.  An \emph{interpolation equation} $E$
is $x\,w_1\dots w_k=u$ with $x:A=(A_1,\dots,A_k,\Ground)$, closed lnf
terms $w_i:A_i$, and a closed lnf goal $u:\Ground$.  A \emph{solution} is
a closed lnf $t:A$ with $t\,w_1\dots w_k=_{\beta\eta}u$.  The order of
$E$ is $\ord(A)$ and its arity is $\ar(A)$.
\end{definition}

Under (R2) every constant of $u$ is first order, so $u$ is a first-order
ground term.  $\CE$ is the set of constants of $u$ together with a fresh
$b:\Ground$ [S, Def.~10].  $\Sub(u)$ is the set of subterms of $u$, with
subterm order $\sq$.  $\depth(a)=0$ and
$\depth(f\,u_1\dots u_m)=1+\max_i\depth(u_i)$.  $|u|$ is the number of
nodes and $\br(u)$ the number of leaves (branches) of $u$.

\begin{lemma}[S, Prop.~4, Fact~6]\label{lem:reduces}
\Paper{} (1) If $t$ solves $E$ then $t\,w_1\dots w_k\to_\beta^*u$.
(2) If $E$ has a solution, it has one with all constants in $\CE$.
\end{lemma}

\begin{definition}[Trees; S, \S3.1]
The tree of an lnf term uses dummy lambdas.  $\lambda\vec y.\,s$ with
$s:\Ground$ ($\vec y$ possibly empty) is a node labelled $\lambda\vec y$
with one successor.  $h\,s_1\dots s_m$ is a node labelled $h$ with
successors $s_1,\dots,s_m$.  $\suc ni$ is the $i$-th successor of $n$.
The \emph{binding relation} $n\bnd m$ holds when $n$ is labelled
$\lambda y_1\dots y_p$ and binds the variable occurrence at $m$.  Nodes
are \emph{lambda}, \emph{variable} or \emph{constant} nodes; a lambda
node with non-empty $\vec y$ is a \emph{binding node}.

The \emph{interpolation tree} $@\,t\,w_1\dots w_k$ has a root
$\varepsilon$ labelled $@$ with successors $1,\dots,k+1$ carrying the
trees of $t,w_1,\dots,w_k$.
\end{definition}

%======================================================================
\section{Games and interpolation [S, \S4]}\label{sec:games}
%======================================================================

Throughout, $E$ is $x\,w_1\dots w_k=u$ and $t:A$ is a closed lnf term
(a \emph{potential solution}).

\subsection{The game}

\begin{definition}[States, look-up tables; S, Defs.~16--17]\label{def:tables}
$N$ is the set of nodes of $@\,t\,w_1\dots w_k$ and $N_1\subseteq N$ the
set of binding nodes.  The set of \emph{states} is
$S=\{[r]\mid r\in\Sub(u)\}\cup\{[\forall],[\exists]\}$; $[\forall]$ and
$[\exists]$ are \emph{final}.  Look-up tables are defined in layers:
$\Theta_1=\{\varnothing\}$, and $\Theta_{i+1}$ is the set of partial maps
\[
N_1\rightharpoonup \textstyle\bigcup_p N^p\times\bigcup_{j\le i}\Theta_j .
\]
An entry $\theta(m)=((m_1,\dots,m_p),\theta')$ says that the variables
bound by $m$ stand for the subtrees at $m_1,\dots,m_p$, read under the
table $\theta'$.  We write $\args_\theta(m)=(m_1,\dots,m_p)$ and
$\env_\theta(m)=\theta'$.  Two tables are \emph{equal} if they have the
same domain and equal entries; this is well defined because the inner
tables come from earlier positions (Proposition~\ref{prop:tables}(4)).
\end{definition}

\begin{definition}[Moves; S, Def.~19 and Fig.~5]\label{def:moves}
Let the current position be $n[r]\theta$ with $r$ not final.  The next
position depends on the label of $n$.
\begin{description}
\item[$@$:] $\suc n1[r]\theta'$ with
  $\theta'=\theta\{((\suc n2,\dots,\suc n{(k+1)}),\theta)/\suc n1\}$.
\item[$\lambda\vec y$:] $\suc n1[r]\theta$.
\item[$a:\Ground$:] $n[\exists]\theta$ if $r=a$, and $n[\forall]\theta$
  otherwise.
\item[$f$ of arity $p>0$:] if $r=f\,r_1\dots r_p$, then
  $\forall$ chooses $j\in\{1,\dots,p\}$ and the next position is
  $\suc nj[r_j]\theta$; otherwise it is $n[\forall]\theta$.
\item[$y_j:\Ground$:] if $m\bnd n$ and $\theta(m)=((m_1,\dots,m_l),\theta')$
  then $m_j[r]\theta'$.
\item[$y_j$ of type $(B_1,\dots,B_p,\Ground)$, $p>0$:]\ \\ if $m\bnd n$ and
  $\theta(m)=((m_1,\dots,m_l),\theta')$ then $m_j[r]\theta''$ with
  $\theta''=\theta'\{((\suc n1,\dots,\suc np),\theta)/m_j\}$.
\end{description}
The last two moves are \emph{jumps}.
\end{definition}

\begin{definition}[Plays, winning; S, Defs.~18, 20]
A \emph{play} of $\Gm(t,E)$ is a finite sequence
$\pi=\pi_1\dots\pi_l$ of positions $\pi_i=n_i[r_i]\theta_i$ with
$\pi_1=\varepsilon[u]\varnothing$, each $\pi_{i+1}$ a next position of
$\pi_i$, $r_i$ non-final for $i<l$, and $r_l$ final.  $\forall$
\emph{wins} $\pi$ if $r_l=[\forall]$.  She \emph{wins the game} if she
wins some play, and \emph{loses} it otherwise.  We write
$\pi\in\Gm(t,E)$, $\iv ij$ for $\pi_i\dots\pi_j$, and say that $\pi_i$ is
\emph{at} $n_i$.
\end{definition}

\subsection{Basic properties of plays}

\begin{proposition}[S, Prop.~12]\label{prop:tables}
\Paper{} Let $\pi\in\Gm(t,E)$, $j\le l$, and
$\theta_j(m)=((m_1,\dots,m_p),\theta')$.
\begin{enumerate}
\item $\theta_j(n)$ is defined iff $n$ is a binding node above $n_j$.
\item $\theta'(n)$ is defined iff $n$ is a binding node strictly above
  every $m_k$.
\item $m$ is labelled $\lambda y_1\dots y_p$, and $y_k$ has the type of
  $m_k$.
\item $\theta'=\theta_i$ for some $i<j$, and either $i=1$ and
  $m_k=\suc{\varepsilon}{(k+1)}$, or $n_i$ is a variable node and
  $m_k=\suc{n_i}k$.
\end{enumerate}
\end{proposition}
\begin{proof}
Simultaneous induction on $j$.  For $j=1$ the table is empty.  For the
step, look at the move from $\pi_{j-1}$.
\begin{itemize}
\item $@$: $j=2$; the only entry is for node $1$, which is above itself;
  its inner table is $\theta_1=\varnothing$ and nothing lies strictly above
  the roots of the $w_i$.
\item $\lambda$, $a$, $f$: the table is unchanged and $n_j$ is not a
  binding node, or is the successor of a constant, which binds nothing
  since constants of $u$ are first order.
\item jump at $y_q$ bound by $m$, with
  $\theta_{j-1}(m)=((m_1,\dots,m_p),\theta')$: then $n_j=m_q$, and
  $\theta_j=\theta'$ or
  $\theta_j=\theta'\{((\suc{n_{j-1}}1,\dots),\theta_{j-1})/m_q\}$.  By the
  induction hypothesis~(2) for $\theta'$, the binding nodes above $m_q$
  are those strictly above $m_q$ (domain of $\theta'$), plus $m_q$
  itself when $y_q$ has higher type.  This gives (1).  Entries of
  $\theta'$ satisfy (2)--(4) by induction.  The new entry has inner
  table $\theta_{j-1}$, an earlier table, and arguments the successors
  of the variable node $n_{j-1}$.  Its types match because $y_q$ and
  $m_q$ have the same type, by induction~(3).  \qedhere
\end{itemize}
\end{proof}

In particular play never gets stuck: at a variable node the binder has
an entry, and a jump from a variable of type $A$ lands at a lambda node
of type $A$.

\begin{definition}[Associated term; S, proof of Thm.~13]\label{def:assoc}
For a node $n$ and a table $\theta$, $\theta[n]$ is the term
\begin{itemize}
\item $\theta[\suc n1]\,\theta[\suc n2]\dots\theta[\suc n{(k+1)}]$ if $n$ is
  labelled $@$;
\item $a$ if $n$ is labelled $a$, and $f\,\theta[\suc n1]\dots\theta[\suc np]$
  if it is labelled $f$;
\item $\theta[\suc n1]$ for a dummy lambda, and
  $\lambda\vec z.\,\theta[\suc n1]$ for $\lambda\vec z$ with $\vec z$
  non-empty;
\item for $y_j$ bound by $m$: $\theta'[m_j]\,\theta[\suc n1]\dots\theta[\suc nq]$
  if $\theta(m)=((m_1,\dots),\theta')$, and $y_j\,\theta[\suc n1]\dots$
  if $\theta(m)$ is undefined.
\end{itemize}
\end{definition}

\begin{theorem}[S, Thm.~13]\label{thm:game}
\Paper{} $\forall$ loses $\Gm(t,E)$ iff $t$ solves $E$.
\end{theorem}
\begin{proof}
Use \emph{leftmost reduction sequences} with the rules
$(\lambda y_1\dots y_p.s)\,s_1\dots s_p\to s\{s_1/y_1,\dots\}$ and
$f\,s_1\dots s_p\to s_j$ (reduction under a constant, which commits to a
branch).  A term is \emph{terminal} if neither rule applies.  A sequence
$s_1\to\dots\to s_m$ is \emph{compatible} with a binder-free normal
$v:\Ground$ if:
\begin{itemize}
\item for $m=1$: if $s_1$ is terminal then $s_1=v$;
\item for a $\beta$-step: the rest is compatible with $v$;
\item for a step $f\,\vec s\to s_j$: $v=f\,\vec v$ and the rest is
  compatible with $v_j$.
\end{itemize}
A \emph{padding} inserts steps $s\to s$.

\emph{Claim.} For every reduction sequence $\sigma$ from
$t\,w_1\dots w_k$ compatible with $u$ there are a play prefix $\pi$ and
a padding $s_1\to\dots\to s_m$ of $\sigma$ with
$s_i=\theta_{2i-1}[n_{2i-1}]$ for all $i$.  Conversely, every play prefix
has such a sequence.

The proof is by induction on the length, by cases on the label of
$n_{2s-1}$, which is never a lambda node (Fact~\ref{fact:parity}).
\begin{itemize}
\item $@$: $\theta_3[n_3]$ is the body of $t$ with $w_i$ for $x_i$.
\item $f$: $\forall$ follows the branch of the reduction step.
\item $y_j$ of higher type: $\theta_{2s-1}[n_{2s-1}]=\theta'[m_j]\,\vec s$
  is a redex $(\lambda\vec y.s')\vec s$.  The next odd position
  represents $s'\{\vec s/\vec y\}$, because the new entry for $m_j$ maps
  $\vec y$ to the successors of $n_{2s-1}$ under $\theta_{2s-1}$.
\item $y_j:\Ground$: $\theta_{2s-1}[n_{2s-1}]=\theta'[m_j]$ and the play
  moves on without a reduction step; insert a padding step.  Only
  finitely many paddings happen in a row, since each uses a table of
  an earlier position.
\end{itemize}
By strong normalisation, plays are finite.

If $t$ solves $E$, every maximal reduction sequence is finite and
compatible with $u$.  So every play ends with $[\exists]$, because a
play ending with $[\forall]$ would give an incompatible sequence.
Conversely, if $\forall$ loses, then every branch of $u$ is followed by
some play, so by some compatible sequence, and $t\,\vec w$ normalises to
$u$.
\end{proof}

\begin{fact}[S, Fact~10]\label{fact:branches}
\Paper{} If $t$ solves $E$, then $\Gm(t,E)$ has exactly $\br(u)$ plays.
\end{fact}
\begin{proof}
The game is deterministic except for $\forall$'s choices at constants.
If $t$ solves $E$, the state at an $f$-node is always $f\,\vec r$
(otherwise $\forall$ wins), so every choice is available.  Every play
ends at some $a$ with state $[a]$.  Hence plays correspond bijectively
to the branches of $u$.
\end{proof}

\begin{fact}[S, Fact~11]\label{fact:parity}
\Paper{} Let $\pi\in\Gm(t,E)$ have length $l$.
(1) If $i<l$ is even, $n_i$ is a lambda node.
(2) If $i>1$ is odd, $n_i$ is a variable or constant node.
(3) $n_l$ is a constant node and $r_l$ is final.
(4) If $i<j<l$ then $r_j\sq r_i$.
\end{fact}
\begin{proof}
In an lnf tree, lambda nodes and non-lambda nodes alternate along paths.
The $@$-move, the $f$-move and the jumps lead to lambda nodes.  The
$\lambda$-move leads to a non-lambda node.  The $a$-move ends the play
at the same node.  States change only at constants, to a subterm.
\end{proof}

\begin{proposition}[S, Prop.~14]\label{prop:distinct}
\Paper{} If $i\ne j$ and $n_i=n_j$ then $\theta_i\ne\theta_j$.
\end{proposition}
\begin{proof}
If $n_i=n_j$ and $\theta_i=\theta_j$, the associated terms coincide.
By the claim in the proof of Theorem~\ref{thm:game} there is a reduction
sequence from that term back to itself that is not all padding,
contradicting strong normalisation.
\end{proof}

\begin{lemma}[Predecessors]\label{lem:pred}
\New{} If $1<j<l$ and $n_j$ is not a lambda node, then $n_{j-1}$
is the lambda node directly above $n_j$, and $\theta_{j-1}=\theta_j$,
$r_{j-1}=r_j$.
\end{lemma}
\begin{proof}
Only the $\lambda$-move leads to a non-lambda node without ending the
play.
\end{proof}

\begin{lemma}[States change only at constants]\label{lem:states}
\New{} \Coq{} Let $t$ solve $E$ and $i<l$.  Then $r_{i+1}\ne r_i$ iff
$n_i$ is a constant node of positive arity, and then $r_{i+1}\ssq r_i$.
Hence for $i\le j<l$: $r_j\ssq r_i$ iff some $n_h$, $i\le h<j$, is a
constant node of positive arity.
\end{lemma}
\begin{proof}
Only the $f$- and $a$-moves change states; the $a$-move ends the play.
At an $f$-node, $\forall$ loses, so $r_i=f\,\vec r$ and $r_{i+1}=r_j\ssq
r_i$.  The second claim then follows from Fact~\ref{fact:parity}(4).
\end{proof}

\begin{lemma}[Locality]\label{lem:locality}
\New{} \Coq{} Let $t,t'$ be lnf terms of type $A$ and $X$ a
prefix-closed set of nodes on which their interpolation trees agree
(labels, successors, binders).  If no play of $\Gm(t,E)$ has a position
outside $X$, then $\Gm(t',E)=\Gm(t,E)$.
\end{lemma}
\begin{proof}
A move depends only on the current state and table, and on the label,
successors and binder of the current node.  By induction, every
sequence of moves that stays in $X$ is the same in both games.  A
sequence for $t'$ that left $X$ for the first time would, up to that
point, be a sequence for $t$ leaving $X$.
\end{proof}

\subsection{Parent and child positions [S, \S4.3]}

\begin{definition}[Parent; S, Defs.~21--22]\label{def:parent}
Let $\pi\in\Gm(t,E)$.
\begin{enumerate}
\item If $n_j$ is a variable node, $\pi_i$ is the \emph{parent} of
  $\pi_j$ if $n_i\bnd n_j$ and $\theta_i(n_i)=\theta_j(n_i)$.
\item If $n_j$ is a lambda node and $n$ is the node directly above
  $n_j$: if $n=\varepsilon$, the parent is $\pi_1$; if $n$ is a constant
  node, the parent is $\pi_{j-1}$; if $n$ is a variable node, the parent
  is the $\pi_i$ such that $\pi_{i+1}$ is the parent of $\pi_{j-1}$ in the
  sense of item~1.
\end{enumerate}
$\pi_j$ is a \emph{child} of $\pi_i$ if $\pi_i$ is its parent.  In
item~2, $n_j$ is a successor of $n_i$.
\end{definition}

\begin{proposition}[S, Prop.~15, Fact~16]\label{prop:parent}
\Paper{} Every position $\pi_j$, $j>1$, at a variable node or at a lambda
node has a unique parent $\pi_i$, and $i<j$.
\end{proposition}
\begin{proof}
Variable node $n_j$ with binder $m$: an entry for $m$ enters a table
only when play is at $m$ (the $@$-move or a higher-type jump), so some
earlier $\pi_i$ at $m$ has $\theta_i(m)=\theta_j(m)$.  If there were two,
$\pi_i$ and $\pi_{i'}$, then $\pi_{i-1}$ and $\pi_{i'-1}$ would be at the
same node with the same table (the entry records both), contradicting
Proposition~\ref{prop:distinct}.  Lambda node: use the three cases of
Definition~\ref{def:parent}(2) and the variable case.
\end{proof}

\begin{definition}[Chain; S, Def.~23]\label{def:chain}
For $\pi_j$ at a lambda node, a \emph{chain} for $\pi_j$ is a sequence
$\pi_{j_1},\dots,\pi_{j_{2p}}$ with
(1) $\pi_{j_{2p}}=\pi_j$;
(2) $\pi_{j_{2m-1}}$ is the parent of $\pi_{j_{2m}}$ for $1\le m\le p$;
(3) $j_{2m+1}=j_{2m}+1$ for $1\le m<p$.
The nodes $n_{j_1},\dots,n_{j_{2p}}$ then form a downward path.
\end{definition}

\begin{definition}[S, Def.~24]
$\theta\le\theta'$ ($\theta'$ \emph{extends} $\theta$) if $\theta'(m)=\theta(m)$
whenever $\theta(m)$ is defined.
\end{definition}

\begin{proposition}[S, Props.~17--18]\label{prop:mono}
\Paper{}
(1) If $n_j$ is a lambda node with parent $\pi_i$, then
$\theta_i\le\theta_j$.
(2) For a chain $\pi_{j_1},\dots,\pi_{j_{2p}}$:
$\theta_{j_1}\le\cdots\le\theta_{j_{2p}}$, and if $n_{j_a}\bnd n_{j_b}$
then $\pi_{j_a}$ is the parent of $\pi_{j_b}$.
\end{proposition}
\begin{proof}
(1) Cases on the label of $n_i$.
\begin{itemize}
\item $@$: $\theta_1=\varnothing$.
\item Constant: $\theta_j=\theta_i$.
\item Variable $z$: $\pi_{i+1}$ is the parent of $\pi_{j-1}$ (at $y_k$),
  so $\theta_{j-1}(n_{i+1})=\theta_{i+1}(n_{i+1})=((\dots),\theta_i)$.
  Hence $\theta_j=\theta_i$ or $\theta_i\{\dots/n_j\}$, and $\theta_i$
  has no entry for $n_j$ (Proposition~\ref{prop:tables}(1)).
\end{itemize}
(2) Chain links are parent/child pairs, and the step $j_{2m}\to
j_{2m}+1$ is a $\lambda$-move, which keeps the table.  The second claim
follows with Proposition~\ref{prop:parent}.
\end{proof}

\begin{proposition}[S, Prop.~19]\label{prop:agree}
\Paper{} If $\theta_i(m)=\theta_j(m)$ are defined and $m'$ is a binding
node above $m$, then $\theta_i(m')=\theta_j(m')$.
\end{proposition}
\begin{proof}
Take the chains for $\pi_i$ and $\pi_j$ from the node above $m'$, and
their prefixes ending at $m$.  The last positions of the prefixes have
equal entries for $m$.  So they coincide, since otherwise their
predecessors would be at the same node with the same table
(Proposition~\ref{prop:distinct}).  Hence the prefixes coincide, and
monotonicity gives the claim.
\end{proof}

\begin{lemma}[Ancestors are visited first]\label{lem:ancestors}
\New{} \Coq{} If $\pi_j$ is at $n$ and $m\ne\varepsilon$ is a
non-lambda node strictly above $n$, then some $i<j$ has $n_i=m$.  If $n$
is a lambda node, there is a unique chain for $\pi_j$ whose first
position is at $m$.
\end{lemma}
\begin{proof}
Let $\rho(j)=j-1$ at non-lambda nodes (Lemma~\ref{lem:pred}) and
$\rho(j)=$ the parent index at lambda nodes.  Then $n_{\rho(j)}$ is the
node directly above $n_j$, and iterating $\rho$ walks up the path.
The visited positions form a chain, and $\rho$ is a function, which
gives uniqueness.
\end{proof}

\begin{lemma}[Alternation]\label{lem:alt}
\New{} If $n_j$ is a variable node then $n_{j+1}$ is a lambda node, and
exactly one of $n_j,n_{j+1}$ lies in the $t$-component.
\end{lemma}
\begin{proof}
By Proposition~\ref{prop:tables}(4), every entry for a binder in $w_q$
points into $t$ and vice versa, since every $w_q$ is closed and $t$ is
closed.
\end{proof}

%======================================================================
\section{Tree automata revisited [S, \S5]}\label{sec:automata}
%======================================================================

This section is not used by the decidability proof.  We record the
definitions and the statements; proofs are as in the paper and have not
been re-checked here.

\begin{definition}[Variable profiles, modes; S, Def.~25]
Let $V$ be a finite set of variables.  For $y\in V$ the set $\Gamma(y)$
of \emph{$y$-profiles} is defined by induction on the type of $y$.
\begin{itemize}
\item If $y:\Ground$, then $\Gamma(y)=\{(y,r,\varnothing)\mid r\in\Sub(u)\}$.
\item If $y:(B_1,\dots,B_p,\Ground)$, then
  $\Gamma(y)=\{(y,r,\Gamma)\mid r\in\Sub(u),\
  \Gamma\subseteq\bigcup_{i\le p}\bigcup_{z:B_i\in V}\Gamma(z)\}$.
\end{itemize}
A \emph{mode} is a pair $(r,\Gamma)$ with $r\in\Sub(u)$ and
$\Gamma\subseteq\bigcup_{y\in V}\Gamma(y)$.  A profile abstracts a position
at a $y$-node together with the profiles of the children of the next
position.
\end{definition}

\begin{definition}[The automaton $A_E$; S, Def.~26]
Fix a finite alphabet $\Sigma$ and well-named interpolation
$\Sigma$-trees.  $A_E=(Q,\Sigma,F,\Delta)$ is a bottom-up tree automaton.
Its states are finite sets of modes, and its accepting set is
$F=\{\{(u,\varnothing)\}\}$.  The transitions are given by cases on the
label.
\begin{itemize}
\item $a\Rightarrow\Lambda$ for $\Lambda\subseteq\{(a,\varnothing)\}$.
\item $y:\Ground$: $y\Rightarrow\Lambda$ for
  $\Lambda\subseteq\{(r,\{(y,r,\varnothing)\})\}$.
\item $f\Lambda_1\dots\Lambda_p\Rightarrow\Lambda$ when:
  \begin{itemize}
  \item each $(r,\Gamma)\in\Lambda$ has $r=f\,r_1\dots r_p$, with
    $(r_i,\Gamma_i)\in\Lambda_i$ and $\Gamma=\bigcup\Gamma_i$;
  \item conversely, each $(r_i,\Gamma_i)\in\Lambda_i$ arises in this way.
  \end{itemize}
\item $y\Lambda_1\dots\Lambda_p\Rightarrow\Lambda$ for $y$ of higher type
  when:
  \begin{itemize}
  \item each $(r,\Gamma)\in\Lambda$ contains a profile $(y,r,\Gamma')$
    whose sub-profiles $(z_i,r_{ij},\Gamma'_{ij})$ are matched by modes
    $(r_{ij},\Gamma'_{ij}\cup\Sigma_{ij})\in\Lambda_i$, with
    $\Gamma=\bigcup\Sigma_{ij}\cup\{(y,r,\Gamma')\}$;
  \item conversely, each mode of $\Lambda_i$ arises in this way.
  \end{itemize}
\item $\lambda\vec y\,\Lambda\Rightarrow\Lambda$.
\item $@\Lambda_1\dots\Lambda_{k+1}\Rightarrow\{(u,\varnothing)\}$ when
  $\Lambda_1=\{(u,\Gamma)\}$, every $(z,r',\Gamma')\in\Gamma$ has $z=x_i$
  and $(r',\Gamma')\in\Lambda_{i+1}$, and conversely.
\end{itemize}
\end{definition}

\begin{theorem}[S, Thm.~20]
\Paper{} For a fixed finite alphabet and a well-named $\Sigma$-tree,
$A_E$ accepts $@\,t\,w_1\dots w_k$ iff $t$ solves $E$.
\end{theorem}

\begin{definition}[Dependency tree automata; S, Defs.~27--28]
$A=(Q,\Sigma,q_0,\Delta)$ with $\Sigma=\Sigma_1\cup\Sigma_2\cup\Sigma_3$
(binders, bound variables, other symbols) is a top-down automaton with
rules of three forms:
\begin{enumerate}
\item $qs\Rightarrow(q_1,\dots,q_k)$ for $s\in\Sigma_2\cup\Sigma_3$ of arity
  $k$;
\item $qs\Rightarrow q's'$ for $s\in\Sigma_1$, $s'\in\Sigma_3$;
\item $(q',q)s\Rightarrow q_1x$ for $s\in\Sigma_1$, $x\in\Sigma_2$.
\end{enumerate}
In a rule of the third form, $q'$ is the state at the binder of $x$.
A run labels nodes top-down from $(\varepsilon,q_0)$.  It is accepting
if it reaches every leaf.
\end{definition}

\begin{theorem}[S, Thm.~21]
\Paper{} For dependency tree automata, non-emptiness is decidable in
time $O(n^4)$, and the languages are closed under intersection and
union.
\end{theorem}

\begin{definition}[Alternating dependency automata, $A_E$; S, Defs.~29--31]
Rule~1 becomes $qs\Rightarrow(Q_1,\dots,Q_k)$ with $Q_i\subseteq Q$.  A
run spawns one copy at $\suc ni$ for each $q'\in Q_i$.  For $E$, the
states are the modes, $q_0=(u,\varnothing)$, and the rules follow the
game:
\begin{itemize}
\item $@$ chooses the profiles of the $x_i$;
\item a lambda above a variable $z_i$ reads a $z_i$-profile from the
  state at its binder;
\item a variable of positive arity spawns the profiles of its arguments;
\item a constant $f$ splits $f\,r_1\dots r_k$ into the $r_i$.
\end{itemize}
\end{definition}

\begin{theorem}[S, Thm.~22]
\Paper{} The alternating dependency automaton $A_E$ accepts the binding
tree $@\,t\,w_1\dots w_k$ iff $t$ solves $E$.  Non-emptiness of
alternating dependency automata is undecidable in general (Ong and
Tzevelekos), so this does not by itself give decidability.
\end{theorem}

%======================================================================
\section{Uniformity of game playing [S, \S6.1]}\label{sec:uniformity}
%======================================================================

\begin{proposition}[S, Prop.~23]\label{prop:between}
\Paper{} If $j<m$ and $\pi_i$ is the parent of both $\pi_j$ and $\pi_m$,
then $\pi_j$ is the parent of some $\pi_k$ with $j<k<m$.
\end{proposition}
\begin{proof}
\emph{Variable case.} $n_i\bnd n_j,n_m$ and
$\theta_j(n_i)=\theta_m(n_i)=((m_1,\dots,m_p),\theta_l)$.  Then $\pi_{j+1}$ is
at some $m_q$ with $\theta_{j+1}=\theta_l\{((\suc{n_j}1,\dots),\theta_j)/m_q\}$.
$\theta_l$ has no entry for $n_i$, since the jump crosses between $t$ and
some $w$.  An entry of $\theta_l$, or of a table reachable from it,
cannot be the entry of $n_i$ recorded at $\pi_i$, by
Proposition~\ref{prop:distinct}.  So the only way back to a table
containing that entry is through a variable bound by $m_q$, i.e.\
through a child of $\pi_j$.

\emph{Lambda case.} Cases on the node above.  For $@$, play leaves $t$
(resp.\ $w_q$) by a jump from a child of $\pi_j$.  A constant node has
only one child position.  For a variable node, apply the variable case
to $\pi_{i+1}$, $\pi_{j-1}$, $\pi_{m-1}$.
\end{proof}

\begin{definition}[Correspondence; S, Def.~32]
$\iv i{i+m}$ and $\iv j{j+m}$ \emph{correspond} if $n_{i+h}=n_{j+h}$ for
$0\le h\le m$.
\end{definition}

\begin{proposition}[S, Prop.~24]\label{prop:24}
\Paper{} Let $\pi_i$ be the parent of $\pi_j$ and of $\pi_m$, $j<m$, and
let $n_j,n_m$ carry the same variable.  Then there is $k$ such that
$\iv{j+1}{j+k}$ and $\iv{m+1}{m+k}$ correspond and either
\begin{enumerate}
\item[(1)] $\pi_j$ is the parent of $\pi_{j+k+1}$, $\pi_m$ is the parent of
  $\pi_{m+k+1}$, and $n_{j+k+1}=\suc{n_j}p$ iff $n_{m+k+1}=\suc{n_m}p$; or
\item[(2)] $n_{j+k}$ is a constant node, and $\pi_{m+k}$ is final or
  $n_{j+k+1}\ne n_{m+k+1}$.
\end{enumerate}
\end{proposition}
\begin{proof}
Write $\theta_i(n_i)=((m_1,\dots),\theta_l)$ and let $n_j,n_m$ carry $y_q$.
Then $n_{j+1}=n_{m+1}=m_q$ and
\[
\theta_{j+1}=\theta_l\{((\suc{n_j}1,\dots),\theta_j)/m_q\},\qquad
\theta_{m+1}=\theta_l\{((\suc{n_m}1,\dots),\theta_m)/m_q\}.
\]
These differ only in the entry for $m_q$.  Hence the two continuations
make the same moves until play reaches either a variable bound by
$m_q$, which gives case~(1) with the same successor index, or a
constant where $\forall$'s choices or the outcomes differ, which gives
case~(2).  This is the special case of Proposition~\ref{prop:24star} in
which the only difference is at $m_q$.  The hypothesis $n_j\ne n_m$ of
the paper is not needed.
\end{proof}

\subsection{A generalisation without common parent}

\begin{definition}[Difference pairs]\label{def:diff}
\New{} For tables $\theta,\theta'$, let $P(\theta,\theta')$ be the least
set of pairs of tables such that
\begin{itemize}
\item $(\theta,\theta')\in P$;
\item if $(\rho,\rho')\in P$ and $\args_\rho(m)=\args_{\rho'}(m)$, then
  $(\env_\rho(m),\env_{\rho'}(m))\in P$.
\end{itemize}
A \emph{difference} of $(\theta,\theta')$ is a triple $(\rho,\rho',m)$ with
$(\rho,\rho')\in P$ and $\args_\rho(m)\ne\args_{\rho'}(m)$.  We write
$\theta\equiv\theta'$ if there is no difference.  Under the reading of
tables as substitutions (Definition~\ref{def:assoc}), $\theta\equiv\theta'$
means that the two tables substitute the same subtrees, hereditarily.
\end{definition}

\begin{proposition}[Prop.~24\textsuperscript{*}]\label{prop:24star}
\New{} Let $a<b$ and $n_a=n_b$.  Then there is $k\ge0$ such that
$\iv a{a+k}$ and $\iv b{b+k}$ correspond and one of the following
holds.
\begin{enumerate}
\item[(1)] \emph{(binder split)} $n_{a+k}$ is a variable $y_q$ bound by
  $m$, and the entries $\theta_{a+k}(m)$, $\theta_{b+k}(m)$ are
  $\rho(m),\rho'(m)$ for a difference $(\rho,\rho',m)$ of
  $(\theta_a,\theta_b)$.  Then $\pi_{a+k+1}$ and $\pi_{b+k+1}$ are at the
  $q$-th successors of the nodes that created these entries.  They are
  children (Definition~\ref{def:parent}) of the creating positions.
\item[(2)] \emph{($\forall$ split)} $n_{a+k}$ is a constant node, and
  $\pi_{a+k}$ or $\pi_{b+k}$ is final, or $\forall$ takes different
  branches.
\end{enumerate}
In particular, if $\theta_a\equiv\theta_b$ then only case~(2) can occur.
\end{proposition}
\begin{proof}
Let $\hat P$ be the closure of $P(\theta_a,\theta_b)$ under the
following operation: from $(\rho,\rho')\in\hat P$ and a pair
$(\sigma,\sigma')\in\hat P$, form
$(\rho\{((\vec n),\sigma)/m\},\rho'\{((\vec n),\sigma')/m\})$, with the
same $\vec n$ and $m$ on both sides.  By induction on $\hat P$: every
entry pair of a table pair in $\hat P$ either has equal args (and an env
pair in $\hat P$), or is an entry pair of some table pair in $P$.

\emph{Invariant:} as long as the runs correspond,
$(\theta_{a+i},\theta_{b+i})\in\hat P$.  Check it move by move
(Definition~\ref{def:moves}):
\begin{itemize}
\item At $\lambda$ and constant nodes the table is unchanged.  At a
  constant node only the state or $\forall$'s choice can separate the
  runs; this is case~(2).
\item At $@$ (only at $a=b=1$, impossible since $a<b$) or at a
  higher-type jump, the new entry has $\vec n$ the successors of the
  common node and env the current pair, so the invariant is kept.
\item At a variable $y_q$ bound by $m$:
  \begin{itemize}
  \item if $\args(m)$ agree, the next nodes agree, and the new tables
    come from the env pair, which lies in $\hat P$;
  \item if they differ, the entry pair is not a created one, so it is
    an entry pair of a pair in $P$.  This is case~(1).
  \end{itemize}
\end{itemize}
The successor and child claims follow from
Proposition~\ref{prop:tables}(4) and Definition~\ref{def:parent}.
Plays are finite, so $k$ exists.
\end{proof}

\begin{corollary}[24$'$: no common parent]\label{cor:24prime}
\New{} Let $j<m$ and let $n_j,n_m$ carry the same variable $y_q$ of a
binder $b$.  Let the parents of $\pi_j,\pi_m$ be visits $\pi_i,\pi_{i'}$
of $b$ with $\theta_i(b)=((\vec m),\theta_l)$ and
$\theta_{i'}(b)=((\vec m),\theta_{l'})$ (same arguments).
\begin{itemize}
\item If $\theta_l\equiv\theta_{l'}$, the conclusion of
  Proposition~\ref{prop:24} holds.
\item Otherwise there is a third outcome: a binder split
  (Proposition~\ref{prop:24star}(1)) at a difference of
  $(\theta_l,\theta_{l'})$.
\end{itemize}
\end{corollary}
\begin{proof}
Apply Proposition~\ref{prop:24star} to $a=j+1$, $b=m+1$.  The
differences of $(\theta_{j+1},\theta_{m+1})$ are the one at $m_q$ (if
$n_j\ne n_m$) together with those of $(\theta_l,\theta_{l'})$.
\end{proof}

\subsection{Complementary intervals and structure of terms}

\begin{definition}[Complementary; S, Def.~33]\label{def:compl}
Let $j<m$ and let $n_j,n_m$ carry the same variable.  The intervals
$\iv{j+1}{j+k}$ and $\iv{m+1}{m+k}$ are \emph{complementary} if case~(1)
or case~(2) of Proposition~\ref{prop:24} holds for them.  No common
parent is required.
\end{definition}

\begin{definition}[Levels, embedded, families; S, Defs.~34--36]\label{def:struct}
\begin{itemize}
\item A variable node has \emph{level}~$1$ if it is bound by the root
  of $t$ or of some $w_q$, and level $j+1$ if it is bound by a successor
  of a level-$j$ node.
\item A node $m$ is \emph{embedded} if some $n\ne m$ above $m$ carries
  the same variable and has the same binder.
\item $n$ is a \emph{descendant under binding} of $m$ if $n=m$, or some
  successor $m'$ of $m$ binds some $n'$ and $n$ is a descendant under
  binding of $n'$.
\item Two nodes are in the same \emph{family} if they have a common
  ancestor under binding.
\item A node is \emph{end} if it has no descendant under binding other
  than itself.
\end{itemize}
\end{definition}

\begin{lemma}[Levels]\label{lem:levels}
\New{} If $\ord(A)\in\{2p,2p+1\}$ and $t\in T_A(\CE)$, every variable
node of $t$ has a level $\le p$.  A level-$j$ variable has order at
most $\ord(A)-2j+1$.
\end{lemma}
\begin{proof}
Constants are first order, so they bind nothing.  A level-$1$ variable
has type $A_i$ with $\ord(A_i)\le\ord(A)-1$.  A successor of
$y:(B_1,\dots,B_m,\Ground)$ binds variables of order
$\le\ord(B_i)-1\le\ord(y)-2$.  Orders are $\ge1$, so $j\le p$.
\end{proof}

\begin{fact}[S, Fact~25]\label{fact:25}
\Paper{} If $n_i$ is end, at most one position $\pi_j$ with $j>i$ is a
child of $\pi_i$.
\end{fact}
\begin{proof}
Suppose $\pi_j$ and $\pi_m$ ($j<m$) are both children of $\pi_i$.  They
are at successors of $n_i$, and by Definition~\ref{def:parent}(2)
$\pi_{i+1}$ is the parent of both $\pi_{j-1}$ and $\pi_{m-1}$.  By
Proposition~\ref{prop:between}, $\pi_{j-1}$ is the parent of some
$\pi_k$ with $j-1<k<m-1$.  By Definition~\ref{def:parent}(2) again, this
means that $\pi_j$ is the parent of $\pi_{k-1}$.  So the successor $n_j$
of $n_i$ binds the variable at $n_{k-1}$, and $n_i$ is not end.
\end{proof}

\begin{fact}[S, Fact~26]\label{fact:26}
\Paper{} If $n_j,n_m$ ($j<m$) are level-$1$ nodes of $t$ carrying the same
variable, then for some $k$ the intervals $\iv{j+1}{j+k}$ and
$\iv{m+1}{m+k}$ are complementary.
\end{fact}
\begin{proof}
Both are bound by the root of $t$, which is visited only at $\pi_2$.  So
$\pi_2$ is their common parent; apply Proposition~\ref{prop:24}.
\end{proof}

\begin{fact}[S, Fact~27]\label{fact:27}
\Sketch{} Let $n,m$ carry the same variable with the same binder $b$,
let $m$ lie below $\suc np$, and let $n_j=m$.  Then there are $i<j$ and
$k$ such that
\begin{enumerate}
\item $n_i=n$ and $\iv{i+1}{i+k}$, $\iv{j+1}{j+k}$ are complementary;
\item if $n,m$ are end and $r_{i+1}=r_{i+k}$, then $n_{j+k+1}=\suc{n_j}p$.
\end{enumerate}
\end{fact}
\begin{proof}[Proof sketch]
(1) By Lemma~\ref{lem:ancestors} the chain for $\pi_{j-1}$ from the
position at $b$ passes a position $\pi_i$ at $n$.  By
Proposition~\ref{prop:mono} it has the same entry for $b$, hence the
same parent as $\pi_j$.  Apply Proposition~\ref{prop:24}.

(2) The chain passes a child of $\pi_i$ at $\suc np$.  If $n$ is end,
that is the only child (Fact~\ref{fact:25}).  In case~(1) of
Proposition~\ref{prop:24} the return is to that child, so the index is
$p$.  Excluding case~(2) needs the non-contribution hypothesis; the
paper does not spell this out.
\end{proof}

\begin{proposition}[S, Prop.~28]\label{prop:28}
\Paper{} Let $\pi_j$ be a position in $w_q$ and
$\pi_{j_1},\dots,\pi_{j_{2p}}$ a chain for $\pi_j$.  If
$\pi_{j_{2a-1}},\pi_{j_{2b-1}}$ ($a\le b$) are positions in $t$, then
$n_{j_{2a-1}}$ and $n_{j_{2b-1}}$ belong to the same family.
\end{proposition}
\begin{proof}
Both are variable nodes, since their children lie in $w_q$.  Extend the
chain upward to the successor of the root of $w_q$.  The position before
the root of $w_q$ is at a level-$1$ node of $t$.  That node is an
ancestor under binding of every $t$-node on the chain, since each jump
back into $w_q$ is through a variable bound inside the chain path.
\end{proof}

\begin{definition}[Recurrence property; S, Def.~37]\label{def:recurrence}
A chain $\pi_{j_1},\dots,\pi_{j_{2p}}$ has the \emph{recurrence property}
if, for every $i\le p$ with $n_{j_{2i-1}}$ a variable node, there are
$j<j_{2i-1}$ and $k$ such that $\iv{j+1}{j+k}$ and
$\iv{m'+1}{m'+k}$ are complementary, where $m'=j_{2i-1}$ and
$j_{2i}=m'+k+1$.
\end{definition}

\begin{proposition}[S, Prop.~29]\label{prop:29}
\Fixed{} Let $n_i$ be a variable or constant node in some $w_q$.  There
is a largest $r\ge0$ such that the \emph{greedy} chain
$\pi_{j_1},\dots,\pi_{j_{2r}}$ with $j_1=i$ has the recurrence property.
Moreover $r$ is at most the depth of $w_q$.
\end{proposition}
\begin{proof}
The greedy chain is built as follows.  Given a recurrent greedy chain of
length $2r'$, let $\pi_{i'}$ be $\pi_i$ (if $r'=0$) or
$\pi_{j_{2r'}+1}$, and let $\pi_{j'}$ be the \emph{first} later child of
$\pi_{i'}$, if any.
\begin{itemize}
\item At a constant node, $j'=i'+1$, and the chain extends.
\item At a variable node, it extends iff for some earlier position at a
  node with the same variable the excursions $\iv{\cdot+1}{\cdot+k}$ and
  $\iv{i'+1}{i'+k}$, $k=j'-i'-1$, are complementary.
\end{itemize}
Otherwise the construction stops.  The nodes form a downward path in
$w_q$, which bounds $r$.
\end{proof}

\begin{remark}
The paper states only ``there is a largest $p$''.  Stages
(Definition~\ref{def:stages}) need one specific chain.  We take the
greedy one, which is what the paper's proof constructs.  A recurrent
chain that uses a later child is not excluded by Definition~\ref{def:recurrence}
alone, so the choice matters.
\end{remark}

%======================================================================
\section{Two transformations [S, \S6.2]}\label{sec:transformations}
%======================================================================

\begin{definition}[Replacement, size; S, \S6.2]
For a variable or constant node $n$ of $t$ and $s:\Ground$, $t[s/n]$
replaces the subtree at $n$ by $s$.  The size measure is
$\mu(t)=(|t|,\nu(t))$, ordered lexicographically: $|t|$ is the number of
nodes (dummy lambdas included) and $\nu(t)$ the number of variable
nodes.  A solution in $T_A(\CE)$ is \emph{smallest} if it is
$\mu$-minimal.
\end{definition}

\begin{definition}[$T_1$; S, Def.~38]
If $n$ is a variable node or a constant node of positive arity and no
play of $\Gm(t,E)$ is ever at $n$, then $t\,T_1\,t[b/n]$.
\end{definition}

\begin{proposition}[S, Fact~30]\label{prop:T1}
\Paper{} \Coq{} If $t$ solves $E$ and $t\,T_1\,t'$, then $t'$ solves $E$
and $\mu(t')<\mu(t)$.
\end{proposition}
\begin{proof}
By Lemma~\ref{lem:ancestors} no play visits a node below $n$ either.
Let $X$ be the set of nodes not below $n$.  Lemma~\ref{lem:locality}
gives $\Gm(t',E)=\Gm(t,E)$, and Theorem~\ref{thm:game} concludes.  If
the removed subtree has one node, it is a variable $y:\Ground$ and
$\nu$ drops; otherwise $|t|$ drops.
\end{proof}

\begin{lemma}[Disjoint contributions]\label{lem:disjoint}
\New{} \Coq{} Let $t$ solve $E$.  If
$\iv{i_1}{j_1},\dots,\iv{i_s}{j_s}$ are intervals of one play with
$r_{i_h}\ne r_{j_h}$ and $j_h\le i_{h+1}$, then $s\le\depth(u)$.
\end{lemma}
\begin{proof}
$r_{i_1}\sqsupset r_{j_1}\sqsupseteq r_{i_2}\sqsupset\cdots$ is a chain in
$\Sub(u)$.
\end{proof}

\begin{lemma}[Constants on a path]\label{lem:constants}
\New{} Let $t$ solve $E$.  On every path of $t$, at most $\depth(u)$
constant nodes of positive arity are visited by some play.
\end{lemma}
\begin{proof}
Take the lowest such constant $c$ and a play that passes from $c$ to its
successor on the path; such a play exists by Fact~\ref{fact:branches}.
The chain from the root (Lemma~\ref{lem:ancestors}) passes every
constant above it with an $f$-move, and each $f$-move strictly
decreases the state.
\end{proof}

\begin{lemma}[Contributing moves over all plays]\label{lem:global}
\New{} Let $t$ solve $E$ and identify positions of plays with a common
prefix.  There are at most $|u|-\br(u)$ distinct $f$-moves in the tree
of plays.
\end{lemma}
\begin{proof}
An $f$-move at state $[r]$ consumes the internal node of $u$ at which
$r$ occurs.  Here states are read as occurrences, determined by
$\forall$'s choices.  Two plays with the same choices up to that point
coincide up to that point.
\end{proof}

\begin{definition}[End path; S, Def.~39]\label{def:endpath}
\Fixed{} A downward path $\gamma=m_1,\dots,m_{2p}$ ($p\ge1$) of $t$ is
\emph{end} if
\begin{enumerate}
\item $m_1$ is a variable node;
\item for $1<i\le p$, $m_{2i-1}$ is a variable node whose binder lies on
  $\gamma$;
\item every variable node strictly below $m_{2p}$ is bound by a node
  strictly above $m_1$ or \emph{strictly} below $m_{2p}$.
\end{enumerate}
\end{definition}

\begin{remark}
The paper says ``below $m_{2p}$''.  It must be read strictly: if
$m_{2p}$ bound a variable below it, that variable would become free after
$T_2$.
\end{remark}

\begin{definition}[Contribution; S, Def.~40]\label{def:contrib}
\Fixed{} An interval $\iv ij$ with $j<l$ \emph{contributes} if
$r_i\ne r_j$.  The paper requires $n_j$ to be a lambda node; we drop
this, so that a final stage $\iv ij$ contributes iff $\iv i{j-1}$ does,
as [S, Def.~44] needs.
\end{definition}

\begin{definition}[Redundancy; S, Def.~41]\label{def:redundant}
Let $\gamma$ be an end path.  $\gamma$ is \emph{redundant in} $\iv ik$ if
either no $j_1\in[i,k)$ has $n_{j_1}=m_1$, or, for the least such $j_1$,
\begin{enumerate}[(a)]
\item there are $j\le k$ and a chain $\pi_{j_1},\dots,\pi_{j_{2p}}$ for
  $\pi_j$ with $n_{j_h}=m_h$ for all $h$;
\item $\iv{j_1}j$ does not contribute;
\item $\gamma$ is redundant in $\iv jk$.
\end{enumerate}
$\gamma$ is \emph{redundant in} $\Gm(t,E)$ if it is redundant in $\iv1l$
for every play.
\end{definition}

\begin{definition}[$T_2$; S, Def.~42]
If $\gamma$ is an end path redundant in $\Gm(t,E)$ and $t'$ is the
subtree at $\suc{m_{2p}}1$, then $t\,T_2\,t[t'/m_1]$.  By
Definition~\ref{def:endpath}(3) the result is closed, and $|t|$ drops.
\end{definition}

\begin{proposition}[S, Prop.~31]\label{prop:T2}
\Sketch{} If $t$ solves $E$ and $t\,T_2\,t''$, then $t''$ solves $E$.
\end{proposition}
\begin{proof}[Proof (paper's argument, with the open step marked)]
Let $\pi\in\Gm(t,E)$.
\begin{enumerate}
\item \emph{Decomposition (proved).} Iterating
  Definition~\ref{def:redundant} gives
  $\pi=\sigma_1\,\iv{i_1}{k_1}\,\sigma_2\cdots\iv{i_q}{k_q}\,\sigma_{q+1}$,
  where each $\iv{i_r}{k_r}$ runs along $\gamma$ by a chain and does not
  contribute, and no position of a $\sigma_r$ is at $m_1$.
\item \emph{No stray positions.} Every position of a $\sigma_r$ below
  $m_1$ is below $m_{2p}$.  The paper argues via the earliest offender:
  its parent is on a chain along $\gamma$, and the end-path condition
  blocks access to the binders on $\gamma$.
\item \emph{Simulation.} Delete the intervals, relocate
  $\suc{m_{2p}}1\,w\mapsto m_1w$, and remove table entries for binders on
  $\gamma$.  This gives a play of $\Gm(t'',E)$ with the same final state,
  and conversely.
\end{enumerate}
Then $\forall$ loses $\Gm(t'',E)$, and Theorem~\ref{thm:game} gives the
claim.
\end{proof}

\begin{gap}[$T_2$ simulation]\label{gap:T2}
\Open{} Steps~2--3 need a precise simulation lemma.  Tables in
$\sigma_{r+1}$ can hold entries created inside $\iv{i_r}{k_r}$, and one
must show that no later position consults them.  The converse direction
is ``almost identical'' in the paper.
\end{gap}

\begin{proposition}[S, Prop.~32]\label{prop:stacking}
\Fixed{} \Sketch{} Let $t$ be a smallest solution, and let
$m_1,\dots,m_p$ be embedded end nodes of $t$ carrying the same variable
with the same binder, with $m_{i+1}$ below $m_i$.  Then
$p\le\depth(u)+1$.
\end{proposition}
\begin{proof}[Proof sketch]
Let $m=\suc{m_p}l$ where $m_p$ is below $\suc{m_{p-1}}l$.  Then $m_p,m$
is an end path.  If it were redundant, $T_2$ would give a smaller
solution.  So some play has an excursion at $m_p$ that contributes, or
that does not return to $m$.  Fact~\ref{fact:27} gives complementary
earlier excursions at $m_1,\dots,m_{p-1}$, which contribute as well.
They are disjoint, so Lemma~\ref{lem:disjoint} applies.
\end{proof}

\begin{gap}[Stacking]\label{gap:stacking}
\Open{} Not argued in the paper: (i) the complementary intervals are
pairwise disjoint and ordered; (ii) the witnessing excursion at $m_p$
is the one to which Fact~\ref{fact:27} applies.  The argument gives
$p-1$ intervals, hence $p\le\depth(u)+1$ rather than the paper's
$p\le\depth(u)$.
\end{gap}

%======================================================================
\section{Partitioning play [S, \S6.3]}\label{sec:partition}
%======================================================================

Throughout this section $t$ solves $E$ and $\pi\in\Gm(t,E)$ has length
$l$.

\begin{definition}[Stages; S, Def.~43]\label{def:stages}
\Fixed{} Stage~$1$ is $\iv12$.  Let stage $N$ end at $\pi_{i-1}$, a
lambda node of $t$, so $n_i$ is its successor.
\begin{enumerate}
\item If $n_i$ is a constant node, stage $N+1$ is $\iv i{i+1}$.
\item If $n_i$ is a variable node, $n_{i+1}$ is a lambda node of some
  $w_q$ (Lemma~\ref{lem:alt}).  Let $\pi_{j_1},\dots,\pi_{j_{2r}}$,
  $j_1=i+2$, be the greedy recurrent chain
  (Proposition~\ref{prop:29}).  Let $p$ be the least $p'\in\{1,\dots,r\}$
  such that $\pi_{j_{2p'}+2}$ is a child of $\pi_i$, or $p=r$ if there is
  none.  Stage $N+1$ is $\iv i{i+3}$ if $p=0$, and $\iv i{j_{2p}+2}$
  otherwise.  The chain $\pi_{j_1},\dots,\pi_{j_{2p}}$ is
  \emph{associated} with the stage.
\end{enumerate}
\end{definition}

\begin{remark}
The paper's definition requires that ``$\iv{j_1}{j_{2p}}$ does not
contain a child of $\pi_i$''.  The text after it uses the cut ``first
$p$ with $\pi_{j_{2p}+2}$ a child of $\pi_i$'', which we adopt.  The
first member $\pi_{j_1}=\pi_{i+2}$ is at a variable or constant node of
$w_q$, not at a lambda node (the margin note ``$\pi_{j_1}$ jest
lambd\k{a}'' on p.~46 of the annotated copy does not fit).  This follows
from Definition~\ref{def:chain}: odd members of a chain are parents of
lambda positions.
\end{remark}

\begin{fact}[S, Fact~33]\label{fact:33}
\Paper{} The stages are consecutive, non-overlapping intervals covering
$\pi$.  Every non-final stage ends at a lambda node of $t$, and the
partition is unique.
\end{fact}
\begin{proof}
Assume stage $N$ ends at a lambda node of $t$.
\begin{itemize}
\item \emph{Constant case:} the stage ends at a successor of the
  constant, or is final.
\item \emph{Variable case, $p=0$:} then $r=0$.  The node $n_{i+2}$ cannot
  be a constant of positive arity, since that would give a recurrent
  chain of length~$2$.  So $\pi_{i+3}$ is a jump into $t$, or the play
  ends at $a$.
\item \emph{Variable case, $p\ge1$:} either $\pi_{j_{2p}+2}$ is a child of
  $\pi_i$, which lies at a successor of $n_i$ in $t$, or $p=r$.  In the
  latter case $\pi_{j_{2r}+1}$ is a variable node, since a constant would
  extend the chain, so $\pi_{j_{2r}+2}$ jumps into $t$, or the play ends.
\end{itemize}
Uniqueness follows from the uniqueness of the greedy chain.
\end{proof}

\begin{definition}[S, Def.~44]
A stage $\iv ij$ is \emph{final} if $r_j$ is final.  A non-final stage
\emph{contributes} if $\iv ij$ does, and a final one if $\iv i{j-1}$
does.  Equivalently (Lemma~\ref{lem:states}), a stage contributes iff
it contains a position at a constant node of positive arity.
\end{definition}

\begin{corollary}\label{cor:contrib-stages}
\New{} In every play at most $\depth(u)$ stages contribute.  Over the
tree of all plays at most $|u|-\br(u)$ distinct stages contribute.
\end{corollary}
\begin{proof}
Lemmas~\ref{lem:disjoint} and~\ref{lem:global}.
\end{proof}

\begin{proposition}[S, Prop.~34]\label{prop:34}
\Paper{} Let $\iv ij$ be stage $N$.
\begin{enumerate}
\item If $n_i$ is labelled $a:\Ground$, then $N$ is final.
\item If $N$ is not final and $n_i$ is a constant or level-$1$ node,
  then $\pi_j$ is a child of $\pi_i$.
\item If $N$ does not contribute and $n_i$ is embedded, then $\pi_j$ is a
  child of $\pi_i$.
\item If $N$ is not final and $n_i$ is a variable node, then $n_j$ is a
  successor of a node in the same family as $n_i$.
\end{enumerate}
\end{proposition}
\begin{proof}
(1) The $a$-move ends the play.

(2) The constant case is Definition~\ref{def:parent}(2).  In the
level-$1$ case, the entry for the root of $t$ is fixed, so $n_{i+1}$ is
the root of $w_q$.  The constants met on the chain bind nothing.  So the
first variable met is bound by that root, and its jump is to a child of
$\pi_i$; this is the cut in Definition~\ref{def:stages}.

(3) By Fact~\ref{fact:27} the excursion is complementary to an earlier
one at the embedding node, which returned to a child.

(4) Extend the associated chain upward to the root of $w_q$ and use
Proposition~\ref{prop:28}.
\end{proof}

\begin{definition}[Comb, path; S, Def.~45]
Let $\iv ij$ be stage $N$.
\begin{itemize}
\item The \emph{comb} $\comb(N)$ is $n_i$ with its successors.
\item The \emph{path} $\pth(N)$ is empty for constant and $@$ stages.
  For a variable stage with associated chain
  $\pi_{j_1},\dots,\pi_{j_{2p}}$ it is
  $n_{i+1},n_{j_1},\dots,n_{j_{2p}},n_{j-1}$, a downward path in some
  $w_q$.
\end{itemize}
Stages inherit \emph{constant/variable}, \emph{level},
\emph{embedded} and \emph{family} from the root $n_i$ of their comb.
\end{definition}

\begin{definition}[Partition tree; S, Def.~46]\label{def:ptree}
The nodes are the stages of $\pi$.  For each non-final stage $N=\iv ij$
there is an edge $N+1\ra\pi K$, where $K\le N$ is the largest stage
whose comb contains $n_j$ and such that $N\ra\pi^*M$ and $K\ra\pi^*M$
for some $M\le N$.
\end{definition}

\begin{remark}[Degeneracy]\label{rem:ptree}
\Fixed{} Every stage reaches stage~$1$, so the side condition always
holds with $M=1$.  As written, the edge goes to ``the latest stage whose
comb contains $n_j$''.  The same node can be the comb root of several
stages, visited with different tables, so this is not the intended
relation.  We use the following reading.
\end{remark}

\begin{definition}[Partition tree, proposed reading]\label{def:ptree2}
\New{} For a non-final stage $N=\iv ij$, let $\pi_h$ be the parent of
$\pi_j$.  Put $N+1\ra\pi K$, where $K$ is the stage containing $\pi_h$.
Gluing $\comb(N+1)$ below the lambda leaf $n_j$ of $\comb(K)$ along every
edge is intended to rebuild the part of $t$ visited by $\pi$.
\end{definition}

\begin{gap}[Partition tree]\label{gap:ptree}
\Open{} Needed:
\begin{enumerate}[(i)]
\item $\pi_h$ is the first position of $K$, so $n_j\in\comb(K)$;
\item the proof of Theorem~\ref{thm:small} goes through with
  Definition~\ref{def:ptree2};
\item every visited node of $t$ lies in the comb of some stage.
\end{enumerate}
\end{gap}

\begin{fact}[S, Fact~35]\label{fact:35}
\Sketch{} If stage $N=\iv ij$ contributes, then one of the following
holds:
\begin{enumerate}
\item $N$ is a constant stage;
\item $\pth(N)$ contains a constant node;
\item some $n_k$, $i<k<j$, is a constant node of $t$ in the comb of an
  earlier stage of the same family.
\end{enumerate}
The third case uses Gap~\ref{gap:ptree}(iii).
\end{fact}

%======================================================================
\section{Unfolding [S, \S6.4]}\label{sec:unfolding}
%======================================================================

\subsection{The paper's definition}

\begin{definition}[Unfolding; S, Def.~47]\label{def:unfold}
Let $N_1<\dots<N_k$ be stages of $\pi$ with combs $c_i$ and paths
$\gamma_i$, such that
\begin{enumerate}[(a)]
\item they belong to the same family and do not contribute;
\item $c_i\ne c_j$ and $\gamma_i\ne\gamma_j$ for $i<j$;
\item for every $j\le k$ (except $j=k$ if $N_k$ is final) there is
  $i\le j$ with $N_j+1\ra\pi N_i$;
\item for every $1<j\le k$ there is $i<j$ such that the root of
  $\gamma_j$ is a successor of the last node of $\gamma_i$.
\end{enumerate}
The \emph{enlarged comb} is built as follows: $(c_1)=c_1$, and
$(c_1\dots c_{j+1})$ places a fresh copy of $c_{j+1}$ below the copy of
the leaf at which $N_j$ finishes.  Variables of the copy are bound by
the copies of their binders when these lie in some $c_i$.  The
\emph{enlarged tree} $(\gamma_1\dots\gamma_j)$ is the union of the paths.
The new stage $(N_1\dots N_j)$ starts at the root of $(c_1\dots c_j)$
and finishes at one of its leaves.  Edges into $N_i$ given by (c) are
redirected to the new stage.
\end{definition}

\begin{fact}[S, Fact~36]\label{fact:36}
\Open{} For each $j\le k$, the play of $(N_1\dots N_j)$ does not
contribute, and it finishes at node $n$ of $c_i$ in $(c_1\dots c_j)$ iff
$N_j$ finishes at $n$ of $c_i$.
\end{fact}

\begin{remark}
The paper does not define the game of the unfolded term (``we do not
provide the game to save notation'').  It also does not prove that the
unfolded term solves $E$.  Section~\ref{sec:nonlinear} shows that, as
written, Definition~\ref{def:unfold} can produce a non-term.
\end{remark}

\begin{definition}[c-embedded; S, Def.~48]
$(c_1\dots c_k)$ is \emph{c-embedded} if its only free variable
occurrence is the head of $c_1$ and some $\alpha$-equivalent enlarged
comb $(c'_1\dots c'_k)$ lies above it in the partition tree.
\end{definition}

\subsection{Order $\le5$: linear heads}

From now on in this subsection, $\ord(A)\le5$.  We write
$t=\lambda z_1\dots z_k.\,S$ and $w_i=\lambda\vec x.\,H_i$.

\begin{lemma}[Shape at order $\le5$]\label{lem:shape5}
\New{} (1) Every level-$2$ variable of $t$ has type $\Ground\to\cdots\to\Ground$
and is an end node.  (2) In $w_i$, every variable bound by a successor
of an $x$-node has base type.  Hence
\[
S::=z_i\,(\lambda\vec y_1.S)\dots(\lambda\vec y_m.S)\mid y\,S\dots S\mid
f\,S\dots S\mid a .
\]
\end{lemma}
\begin{proof}
By Lemma~\ref{lem:levels}, level-$2$ variables have order $\le2$, so
their arguments have base type.  The $x_h$ have order $\le3$, so their
arguments have order $\le2$ and bind variables of order~$\le1$.
\end{proof}

\begin{lemma}[Non-contributing level-1 stages]\label{lem:lvl1stage}
\New{} Let $\pi_i$ be at a $z_i$-node $n$, starting a non-final stage.
The stage does not contribute iff $H_i=x_j\,B_1\dots B_p$ for some $j$.
In that case the stage is $\iv i{i+3}$ and ends at $\suc nj$.
\end{lemma}
\begin{proof}
$\pi_{i+2}$ is at the head of $H_i$.  A constant head contributes, or
ends the play.  A variable head $x_j$ jumps to $\suc nj$, which is a
child of $\pi_i$.
\end{proof}

\begin{definition}[Linear and non-linear heads]\label{def:linear}
\New{} The head $x_j$ of $w_i=\lambda\vec x.\,x_j B_1\dots B_p$ is
\emph{linear} if $x_j$ does not occur in $B_1,\dots,B_p$, and
\emph{non-linear} otherwise.  Repetitions of other variables $x_h$,
$h\ne j$, do not matter here.
\end{definition}

For a $z_i$-node $n=z_i A_1\dots A_m$ of $t$ with $A_j=\lambda\vec y.\,S_j$,
let $ny=y_l R_1\dots R_s$ be an occurrence in $S_j$ with $y_l$ bound at
$\suc nj$.  The \emph{one-step unfolding} is
\[
t'=t\bigl[ny:=z_i\,A_1\dots A_{j-1}\,(\lambda\vec y'.\,y'_l\,\vec R')\,
A_{j+1}\dots A_m\bigr],
\]
with the $A_h$ copied and $\vec R'$ rebound to $\lambda\vec y'$.  This is
Definition~\ref{def:unfold} with $k=2$.

\begin{theorem}[Linear unfolding is sound]\label{thm:linunfold}
\New{} \Coq{} If $x_j$ is a linear head, then $t\,\vec w$ and
$t'\,\vec w$ have a common $\beta$-reduct.  Hence $t'$ solves $E$ iff
$t$ does.
\end{theorem}
\begin{proof}
\begin{enumerate}
\item Contract $w_i\vec A$ at $n$: it reduces to
  $S_j\{\vec y:=\vec B\{\vec x:=\vec A\}\}$.  By linearity, $S_j$
  occurs once.
\item At $ny$, the reduct of $t$ holds
  $B_l\{\vec x:=\vec A\}\,\vec R\{\vec y:=\vec B\{\vec A\}\}$.
\item The reduct of $t'$ holds $w_i\vec A''$, with $\vec A''$ equal to
  $\vec A$ off~$j$.  This reduces to the same term, because
  $\vec B\{\vec x:=\vec A''\}=\vec B\{\vec x:=\vec A\}$ by linearity.
\end{enumerate}
Confluence finishes the proof.  In Coq this is \texttt{unfold\_sound}
(\texttt{Unfolding.v}), proved on raw de~Bruijn terms using
\texttt{confluence} (\texttt{Confluence.v}), with no axioms.  The lnf
version is still to be done.
\end{proof}

\subsection{Non-linear heads: the obstruction and its repair}\label{sec:nonlinear}

\begin{example}[Definition~47 yields no term]\label{ex:obstruction}
\Fixed{} Let $E$ be $X\,w=f\,(f\,a\,d)\,d$ with
$w=\lambda x.\,x\,(\lambda v.\,x\,(\lambda v'.\,v))$ and
$t=\lambda z.\,z\,(\lambda y.\,f\,(y\,a)\,d)$.  The head $x$ is
non-linear.  Number the nodes:
\begin{itemize}
\item in $t$: (1)$\lambda z$, (2)$z$, (3)$\lambda y$, (4)$f$, (5)$\lambda$,
  (6)$y$, (7)$\lambda$, (8)$a$, (9)$\lambda$, (10)$d$;
\item in $w$: (11)$\lambda x$, (12)$x$, (13)$\lambda v$, (14)$x$,
  (15)$\lambda v'$, (16)$v$.
\end{itemize}
In play~B, $\forall$ chooses branch~$1$ at the root $f$ and branch~$2$
at the re-entry:
\[
(0)\,(1)\,(2)\,(11)\,(12)\,(3)\,(4)\,(5)\,(6)\,(13)\,(14)\,(3)\,(4)\,(9)\,(10)\,(10)[\exists].
\]
Its stages under Definition~\ref{def:stages} are
$S_1=\iv12$, $S_2=\iv36$ ($z$, ending at (3)), $S_3=\iv78$ ($f$),
$S_4=\iv9{12}$ ($y$; the chain from $\pi_{11}$ is empty, so it ends at
the re-entry into (3)), $S_5=\iv{13}{14}$ and $S_6=\iv{15}{16}$.

Unfolding $S_2$ into $S_4$ satisfies (a)--(d) of
Definition~\ref{def:unfold}.  In $(c_2c_4)$ the copy of (3) carries the
copy of the $y$-node.  The redirected edge $S_5\to S_2$ must hang the
comb (4) under the same copy of~(3).  So the result is not a term.
\end{example}

\begin{definition}[Tree-wide non-contribution]\label{def:47a}
\Fixed{} Identify stages of different plays that have the same prefix.
In Definition~\ref{def:unfold}(a), read ``do not contribute'' as ``do
not contribute in \emph{any} play passing through the stage''.  This is
also how Stirling's LMCS~2009 version works (Defs.~9.1, 9.5(3)).
\end{definition}

In Example~\ref{ex:obstruction}, plays~A (branch~1 at the re-entry) and
B coincide up to $\pi_{11}$, so they share $S_4$.  In~A the chain
$\pi_{11},\pi_{16}$ is recurrent: $\iv69$ and $\iv{12}{15}$ are
complementary, with common parent $\pi_4$.  So $S_4=\iv9{18}$ passes
the constant at~(4) and contributes.  Under Definition~\ref{def:47a},
$S_2$ is therefore not unfoldable.

\begin{proposition}[Re-entries contribute somewhere]\label{prop:reentry}
\New{} Let $\ord(A)\le5$ and use Definitions~\ref{def:stages}
and~\ref{def:47a}.  Let $N_1=\iv{i_1}{i_1+3}$ be a $z$-stage at $n$ with
head $x_j$ at $\pi_{s-1}$, ending at $\pi_s$ at $\suc nj$.  Let $N_2$
start at $\pi_i$ at a $y_l$-node, and suppose its path starts at a
successor of the head $x_j$ node (Definition~\ref{def:unfold}(d)).
Suppose that in some play~B, $N_2$ ends without contributing at a
re-entry $\pi_e=\pi_{x+1}$ into $\suc nj$ from an inner $x_j$-node at
$\pi_x$.  Assume also the \emph{same-visit} condition: $\pi_{s-1}$ and
$\pi_x$ have the same parent (the visit of $\lambda\vec x$ at
$\pi_{i_1+1}$).  Then $N_2$ contributes in some play~A that shares $N_2$
with~B.
\end{proposition}
\begin{proof}
\emph{Step 1.} By Proposition~\ref{prop:24} for $\pi_{s-1},\pi_x$, there
is $k$ with $\iv s{s-1+k}$ and $\iv{x+1}{x+k}$ corresponding, and
either (1) both runs return to children of $\pi_{s-1}$, resp.\ $\pi_x$,
with the same index (a $\vec y$-node jumps into $B_m$, resp.\ $C_m$), or
(2) they part at a constant~$c$.

\emph{Step 2: case (1) is impossible in B.}  In case~(1) the intervals
are complementary (Definition~\ref{def:compl}).  So the link
$(\pi_x,\pi_{x+k+1})$ is recurrent.  It is also the first child of
$\pi_x$, since successors of $n_x$ are reached only through $\vec y$.
No child of $\pi_i$ occurs in $\iv x{x+k+1}$, since such a child needs a
variable bound in $B_l$, and these occur only inside $\vec C$.  So the
greedy chain of $N_2$ extends past $\pi_x$, and $N_2$ does not end
at~$\pi_e$, a contradiction.

\emph{Step 3: case (2) gives A.}  ``Final'' at $c$ is impossible: the
earlier run would then be final at $\pi_{s-1+k}$ with $s-1+k\le i$, but
it went on to reach the node of $N_2$.  So $\forall$ chose a different
branch at $c$, and $c$ has positive arity.  Let A agree with B up to
$\pi_{x+k}$ and copy the earlier choices from then on; such a play exists
by Fact~\ref{fact:branches}.  By Proposition~\ref{prop:24} again, A
reaches case~(1), at the latest at the copy of the $y_l$-node itself.
So A's chain extends past $\pi_x$, and A's $N_2$ contains $c$.  It
therefore contributes (Lemma~\ref{lem:states}).  A and B coincide up to
$\pi_{x+k}$, beyond $\pi_i$, so they share $N_2$.
\end{proof}

\begin{corollary}[Cross-visit case]\label{cor:crossvisit}
\New{} Drop the same-visit condition: two stages end at the same comb
leaf $L=\suc{n_r}h$ as children of different visits $V_1,V_2$ of $n_r$.
Compare the runs from the two entries into $L$ by
Proposition~\ref{prop:24star}.
\begin{itemize}
\item A split at $L$'s own binder gives complementary intervals
  (Corollary~\ref{cor:24prime}), and the argument of
  Proposition~\ref{prop:reentry} goes through.
\item A $\forall$ split is handled as in Step~3.
\item The new case is a split at an outer binder.  At order~$\le5$ the
  root $z$'s have a single visit, so the only outer binders that can
  differ are the $\lambda\vec y'$ of level-$2$ variables bound above
  $n_r$.
\end{itemize}
Hence, if no $\lambda\vec y'$ lies on the path from the root of $t$ to
$n_r$, the contexts of $V_1,V_2$ are $\equiv$ and the conclusion of
Proposition~\ref{prop:reentry} holds.
\end{corollary}

\begin{gap}[Unfolding]\label{gap:unfold}
\Open{}
\begin{enumerate}[(i)]
\item A precise game for the unfolded term, and Fact~\ref{fact:36}.
\item Soundness of unfolding in general.  It is proved here for linear
  heads (Theorem~\ref{thm:linunfold}).  For non-linear heads, with
  Definition~\ref{def:47a}, the conflict of
  Example~\ref{ex:obstruction} is excluded in the same-visit case
  (Proposition~\ref{prop:reentry}).
\item The remaining cross-visit case of Corollary~\ref{cor:crossvisit}:
  $n_r$ lies below some $\lambda\vec y'$ that its two visits bind
  differently, and play reaches that $\vec y'$ before $\forall$ splits.
\item Consistency across plays: Definition~\ref{def:unfold} composes
  combs per node, not per visit.
\end{enumerate}
\end{gap}

\begin{remark}[Alternatives considered]
\begin{itemize}
\item \emph{Def.~43$'$.} A stage absorbs re-entries into already visited
  $\lambda$-nodes and ends only at (E1) a child of $\pi_i$, (E2) a
  first entry into a new $\lambda$-node, (E3) a $\forall$ split inside
  an absorbed interval, or (F) a final state.  This follows LMCS
  Def.~8.13 (\texttt{order5-unfolding.md}~\S9).  It also blocks
  Example~\ref{ex:obstruction}, but it requires re-proving
  Fact~\ref{fact:33}, Proposition~\ref{prop:34} and Fact~\ref{fact:36}.
  Definition~\ref{def:47a} achieves the same in the same-visit case
  without changing stages.
\item \emph{Exclusion} (LMCS, proof of Thm.~9.11).  Forbid unfolding of
  the tile on which two conflicting tiles depend, and repeat.  On
  Example~\ref{ex:obstruction} this forbids unfolding $S_2$.  The
  paper's count of excluded tiles, $(\delta+p)N(n)$, does not cover this
  tile, which is an ancestor in the family rather than an x-special
  tile or one of its dependents.
\end{itemize}
\end{remark}

%======================================================================
\section{The small solution theorem and decidability}\label{sec:thm37}
%======================================================================

\begin{theorem}[S, Thm.~37]\label{thm:small}
\Open{} Let $E$ have order $2p$ or $2p+1$, and let $q$ bound the arity
of $E$ \emph{and of the constants of $u$}.  If $E$ has a (canonical)
solution, it has one of depth at most $4p^2q^{2p-2}|u|$.
\end{theorem}

\begin{remark}
\Fixed{} The paper lets $q$ be the arity of $E$ only.  Constant nodes of
$t$ branch according to the constants of $u$, so $q$ must bound their
arities too.  For $p=2$ (order~$\le5$) the bound is $16q^2|u|$, and for
[S, Example~12] ($q=2$, $|u|=2$) it is $128$.
\end{remark}

\subsection{Proof architecture}

Let $t$ be a smallest solution.  The paper's argument has the following
steps; ``O'' marks simple observations.

\begin{enumerate}[{O}1.]
\item \Paper{} Nodes of $t$ have level $\le p$ (Lemma~\ref{lem:levels}).
  Branching and binder length are at most~$q$.
\item \Paper{} If a binder binds $(k+1)q$ nodes on one path, then at
  least $k$ of them are embedded nodes carrying the same variable.  The
  binder binds $\le q$ variables, so one variable occurs $\ge k+1$
  times, and all occurrences but the highest are embedded.
\item \Paper{} Level-$p$ nodes are end.
\item \Paper{} There are $\br(u)\le|u|$ plays (Fact~\ref{fact:branches}).
\item \Paper{} $T_1$ and $T_2$ do not apply to $t$.  Hence every variable
  node and every constant node of positive arity is visited, and on each
  path at most $\depth(u)$ constants of positive arity occur
  (Lemma~\ref{lem:constants}).  This depends on Gap~\ref{gap:T2}.
\end{enumerate}

\paragraph{Construction.}
\begin{enumerate}[(a)]
\item Build the partition trees of all plays, glued along common
  prefixes.
\item Unfold exhaustively, starting from stages with no contributing
  descendant under binding (level-$1$ or embedded stages).
\item Apply $T_1$ and $T_2$ exhaustively.
\end{enumerate}
The small term is rebuilt by gluing the (enlarged) combs along the
edges.  A stage $N$ is a \emph{separator} if two plays share all combs up
to and including $N$ but finish $N$ at different nodes.  Separators must
be kept, so that a node entered from several later stages receives one
subterm.

\paragraph{Pigeonhole (S, p.~53).}
\Sketch{} Fix $k$ and take the least $j$ such that some path of $t$ has
$\ge(k+1)q^j$ level-$j$ nodes.
\begin{itemize}
\item If $k$ of them are embedded end nodes with the same variable, then
  $k\le\depth(u)+1$ (Proposition~\ref{prop:stacking}).
\item If they are not end and $k>q$, then some $N_1$ unfolds into $N_2$
  with $(N_1N_2)$ c-embedded.  The reconstruction: there are $\le q$
  ``shapes'' (which successor and which variable leads to $N_2$), so two
  repetitions have the same shape.
\item If $k>q^l$, then $N_1$ unfolds into a c-embedded $(N_2\dots N_r)$
  with $r\le l+1$.
\end{itemize}
Assuming that no $N_j$ in an unfolding is embedded because of an
earlier $N_i$, an enlarged comb has depth $\le2pq^{p-1}$: the layers
contribute $1+q+\dots+q^{p-1}\le pq^{p-1}$, times~$2$ for the lambda
nodes.

\paragraph{Counting (S, pp.~53--54).}
\Sketch{} After unfolding, the stages that $T_2$ cannot remove are:
\begin{enumerate}
\item contributing stages: $\le|u|$ (Corollary~\ref{cor:contrib-stages}),
  or $\le|u|\cdot\br(u)\le|u|^2$ under the tree-wide reading of
  Definition~\ref{def:47a};
\item later non-embedded stages starting from a descendant under
  binding of a contributing comb: $\le pq^{p-1}$ per contributing comb
  (not argued in the paper);
\item separators: $\le pq^{p-1}|u|$ (not argued in the paper).
\end{enumerate}
A path therefore has at most $2pq^{p-1}|u|$ such stages, each with
depth $\le2pq^{p-1}$, which gives $4p^2q^{2p-2}|u|$.

\begin{gap}[Counting]\label{gap:count}
\Open{}
\begin{enumerate}[(i)]
\item Termination of the interleaving ``unfolding grows the term,
  $T_1,T_2$ shrink it''.
\item ``$k>q$ gives a c-embedded $(N_1N_2)$'': the shape pigeonhole is
  not written in the paper.
\item The least-$j$ argument: level-$j$ nodes on one path can be bound
  by different binders.
\item The bounds $pq^{p-1}$ (item~2) and $pq^{p-1}|u|$ (item~3).
\item That $T_2$ applies to \emph{all} remaining stages.  Redundancy is
  defined over all plays, but partition trees are per play.
\item Under Definition~\ref{def:47a}, the factor $|u|$ in the bound
  becomes $|u|^2$.
\end{enumerate}
\end{gap}

\begin{corollary}[Decidability]\label{cor:decidable}
\Paper{} \Coq{} Assume Theorem~\ref{thm:small}.  Then higher-order
matching under (R1), (R2) is decidable.
\end{corollary}
\begin{proof}
A matching equation reduces to the canonical solvability of an
interpolation equation [S, Prop.~5], with constants in $\CE$
(Lemma~\ref{lem:reduces}).  Up to $\alpha$, there are finitely many lnf
terms of bounded depth over $\CE$, since names can be drawn from a
fixed finite set.  Enumerate them and check each by normalisation.  In
Coq this is \texttt{decidable\_interpolation} (\texttt{Small.v}), with
\texttt{Thm37} as a hypothesis.
\end{proof}

%======================================================================
\section{Removing the restrictions [S, \S6.5]}\label{sec:restrictions}
%======================================================================

Summary only; not re-checked.
\begin{itemize}
\item \emph{Several base types.} Nodes carry a result type, which can
  change only at constants, so at most $\depth(u)$ times per play.
  $T_1$ uses one constant $b_B$ per base type, and $T_2$ and unfolding
  are restricted to paths of constant result type.  The bound gets
  $|u|^2$ in place of $|u|$.
\item \emph{Goals with bound variables.} Tables get entries for binders
  below constants, mapping their variables to fresh constants.  Theorem~13
  must be re-proved for the amended game.
\end{itemize}

%======================================================================
\section{Summary of status}\label{sec:summary}
%======================================================================

\begin{center}
\small
\begin{tabular}{lll}
Here & Paper & Status\\\hline
Prop.~\ref{prop:tables}, Thm.~\ref{thm:game}, Prop.~\ref{prop:distinct} & Prop.~12, Thm.~13, Prop.~14 & proved (paper); Coq axioms\\
Props.~\ref{prop:parent}--\ref{prop:agree} & Prop.~15--19 & proved (paper); Coq axioms\\
Lemmas~\ref{lem:pred}--\ref{lem:alt} & implicit & proved; partly Coq\\
Prop.~\ref{prop:between}, \ref{prop:24} & Props.~23--24 & proved (paper)\\
Prop.~\ref{prop:24star}, Cor.~\ref{cor:24prime} & --- & new, proved\\
Facts~\ref{fact:25}, \ref{fact:26} & Facts~25--26 & proved\\
Fact~\ref{fact:27} & Fact~27 & sketch\\
Prop.~\ref{prop:29} & Prop.~29 & corrected (greedy chain)\\
Prop.~\ref{prop:T1} & Fact~30 & proved; Coq\\
Prop.~\ref{prop:T2} & Prop.~31 & sketch, Gap~\ref{gap:T2}\\
Prop.~\ref{prop:stacking} & Prop.~32 & corrected bound, Gap~\ref{gap:stacking}\\
Fact~\ref{fact:33}, Prop.~\ref{prop:34} & Fact~33, Prop.~34 & proved (Coq: uniqueness)\\
Def.~\ref{def:ptree2} & Def.~46 & corrected reading, Gap~\ref{gap:ptree}\\
Thm.~\ref{thm:linunfold} & --- & new; Coq (raw terms)\\
Ex.~\ref{ex:obstruction} & Def.~47 & counterexample\\
Prop.~\ref{prop:reentry}, Cor.~\ref{cor:crossvisit} & --- & new; partial, Gap~\ref{gap:unfold}\\
Fact~\ref{fact:36} & Fact~36 & open\\
Thm.~\ref{thm:small} & Thm.~37 & open, Gaps~\ref{gap:T2}--\ref{gap:count}\\
Cor.~\ref{cor:decidable} & \S6 & proved from Thm.~37; Coq
\end{tabular}
\end{center}

\paragraph{Corrections adopted.}
\begin{enumerate}
\item End paths: ``below $m_{2p}$'' read strictly
  (Definition~\ref{def:endpath}).
\item Contribution of intervals not ending at a lambda node
  (Definition~\ref{def:contrib}).
\item Prop.~29 and stages use the greedy recurrent chain
  (Proposition~\ref{prop:29}).
\item Stage cut: the first $p$ with $\pi_{j_{2p}+2}$ a child of $\pi_i$
  (Definition~\ref{def:stages}).
\item Partition-tree edges via parents (Definition~\ref{def:ptree2}).
\item Prop.~32 gives $p\le\depth(u)+1$.
\item Def.~47(a) read over all plays (Definition~\ref{def:47a}).  As
  written, Def.~47 can produce a non-term
  (Example~\ref{ex:obstruction}).
\item Thm.~37: $q$ also bounds the arities of the constants of $u$.
\end{enumerate}

\paragraph{Open, in order of importance.}
Gap~\ref{gap:unfold} (unfolding: game, soundness, the remaining
cross-visit case), Gap~\ref{gap:count} (counting), Gap~\ref{gap:T2}
($T_2$ simulation), Gap~\ref{gap:ptree} (partition tree),
Gap~\ref{gap:stacking} (stacking).

\end{document}
